
\documentclass[11pt,a4paper]{article}

% ── Packages ──────────────────────────────────────────────────────────
\usepackage[utf8]{inputenc}
\usepackage[T1]{fontenc}
\usepackage{amsmath,amssymb,amsthm,mathtools}
\usepackage{booktabs}
\usepackage{enumitem}
\usepackage{geometry}
\usepackage{hyperref}
\usepackage{xcolor}
\usepackage{tikz}
\usetikzlibrary{arrows.meta,positioning,decorations.pathmorphing,calc}
\usepackage{cleveref}
\usepackage{algorithm}
\usepackage{algorithmic}
\usepackage{microtype}
\usepackage{lmodern}
\usepackage{graphicx}

% ── Geometry ──────────────────────────────────────────────────────────
\geometry{margin=2.5cm}

% ── Hyperref style ────────────────────────────────────────────────────
\hypersetup{
    colorlinks  = true,
    linkcolor   = blue!70!black,
    citecolor   = blue!70!black,
    urlcolor    = blue!60!black,
}

% ── Theorem environments ──────────────────────────────────────────────
\newtheorem{proposition}{Proposition}
\newtheorem{remark}{Remark}
\newtheorem{definition}{Definition}

% ── Custom commands ───────────────────────────────────────────────────
\newcommand{\R}{\mathcal{R}}
\newcommand{\G}{\mathcal{G}}
\newcommand{\KL}{\mathrm{KL}}
\newcommand{\Ber}{\mathrm{Ber}}
\newcommand{\Bern}{\mathrm{Bern}}
\newcommand{\E}{\mathbb{E}}
\newcommand{\I}{I}
\newcommand{\loss}{\mathcal{L}}
\DeclareMathOperator*{\argmin}{arg\,min}
\newcommand{\BCE}{\mathrm{BCE}}

% ══════════════════════════════════════════════════════════════════════
\title{\vspace{-1.0cm}
\textbf{Possible Improvements for the\\Temporal Graph Information Bottleneck}\\[0.4em]
\normalsize\textit{A Mathematical Analysis of the TGIB Objective and Implementation}}
\date{}
\author{%
  Project Report — CTGIB Repository\\[0.2em]
  \small Based on \textit{Self-Explainable Temporal Graph Networks based on\\
  Graph Information Bottleneck} (Seo et al., KDD 2024)}
% ══════════════════════════════════════════════════════════════════════

\begin{document}
\maketitle

% ── Abstract ──────────────────────────────────────────────────────────
\begin{abstract}
\noindent
The Temporal Graph Information Bottleneck (TGIB) framework simultaneously
performs link prediction and explanation generation on temporal graphs by
optimizing a Graph Information Bottleneck (GIB) objective. This report
provides a detailed mathematical analysis of the TGIB formulation and
its open-source implementation, identifying five concrete improvements
that are both theoretically grounded and directly implementable. We
formalise: (1)~an explicit Lagrange multiplier~$\beta$ for the
rate--distortion tradeoff, (2)~a numerically stable closed-form
$\KL$-divergence computation, (3)~temperature annealing for the
Gumbel--Softmax mask sampler, (4)~a configurable and generalisable
prior schedule, and (5)~a $\beta$-warmup strategy to stabilise
optimisation. We provide the complete mathematical derivation for each
improvement, discuss the expected effect on both link-prediction
accuracy and explanation quality, and supply ready-to-use
implementation sketches compatible with the official TGIB codebase.
\end{abstract}

\tableofcontents
\newpage

% ══════════════════════════════════════════════════════════════════════
\section{Introduction}
\label{sec:introduction}
% ══════════════════════════════════════════════════════════════════════

Temporal Graph Neural Networks~(TGNNs) have demonstrated strong
performance in modelling dynamic interactions across social networks,
collaboration graphs, and communication logs. A growing body of work
has therefore focused on \emph{explaining} TGNN predictions---that is,
identifying which past events causally influenced a predicted link
occurrence.

Seo et~al.~\cite{seo2024tgib} proposed \textbf{TGIB}, a \emph{built-in}
explanation framework that unifies prediction and explanation within a
single forward pass. Rather than relying on post-hoc perturbation
methods, TGIB injects stochasticity into the representation of each
past candidate event via learned Bernoulli masks, optimising a
\emph{Graph Information Bottleneck} (GIB) objective. The mask
probabilities naturally serve as importance scores, and the resulting
sparse subgraph constitutes the explanation.

While elegant in formulation, the official implementation contains
several mathematical simplifications and numerical shortcuts that
diverge from the published equations. This report systematically
catalogues these discrepancies and proposes rigorous, drop-in
improvements. Specifically, we address:

\begin{enumerate}[leftmargin=*,itemsep=2pt]
  \item The implicit weight of the rate term ($\beta = 1$).
  \item An ad-hoc $\KL$-divergence computation with an
        epsilon-in-ratio pattern.
  \item A fixed Concrete temperature ($\tau = 1$) with no annealing.
  \item Hardcoded prior parameters ($r_0 = 0.9$, $r_{\min} = 0.7$).
  \item No warmup for the $\beta$ schedule.
\end{enumerate}

Each improvement is derived from first principles, benchmarked against
the paper's formulation, and accompanied by pseudocode for direct
integration.

% ══════════════════════════════════════════════════════════════════════
\section{Background: The TGIB Objective}
\label{sec:background}
% ══════════════════════════════════════════════════════════════════════

\subsection{Information Bottleneck for Temporal Graphs}

Let $e_k = (u_k, v_k, t_k, \mathrm{att}_k)$ denote the $k$-th event in
a temporal graph, and let $\G^k$ denote the $L$-hop computation graph
of~$e_k$ constructed from all events prior to timestamp~$t_k$. TGIB
seeks a \emph{bottleneck code} $\R^k \subseteq \G^k$ that is
simultaneously \emph{sufficient} for predicting the label~$Y_k$
(whether $e_k$ occurs) and \emph{minimal} relative to the full
neighbourhood. Following the GIB principle~\cite{wu2020gib}, this is
formalised as:
\begin{equation}
\label{eq:gib_objective}
\min_{\R^k} \; \underbrace{-\I(Y_k ; \R^k)}_{\text{Distortion (Eq.~\ref{eq:distortion})}}
  \;+\; \beta \underbrace{\I(\R^k ; e_k, \G^k)}_{\text{Rate (Eq.~\ref{eq:rate})}},
\end{equation}
where $\beta > 0$ is a Lagrange multiplier controlling the
accuracy--explanation tradeoff.

\subsection{Variational Upper Bounds}

Since mutual information is intractable to compute exactly, TGIB
optimises variational upper bounds on each term.

\paragraph{Distortion term.}
Using a variational approximation $q_\theta(Y_k \mid \R^k)$ of
$p(Y_k \mid \R^k)$:
\begin{equation}
\label{eq:distortion}
-\I(Y_k ; \R^k) \;\leq\; \E_{Y_k, \R^k}\!\big[-\log q_\theta(Y_k \mid \R^k)\big]
  \;\eqqcolon\; \loss_{\text{cls}}.
\end{equation}

\paragraph{Rate term.}
Decomposing the bottleneck code into per-event Bernoulli masks
$\alpha_j^k \sim \Ber(p_j^k)$ with
$p_j^k = \sigma\!\big(g(X_{e_k}, Z_{e_j})\big)$, and using a
variational marginal $q(\R^k)$ that factorises as a product of i.i.d.
$\Bern(r)$ distributions:
\begin{equation}
\label{eq:rate}
\I(\R^k ; e_k, \G^k)
  \;\leq\; \E_{\G^k}\!\left[\sum_{e_j \in \G^k}
    p_j^k \log \frac{p_j^k}{r}
    + (1 - p_j^k) \log \frac{1 - p_j^k}{1 - r}
  \right]
  \;\eqqcolon\; \loss_{\text{MI}},
\end{equation}
which is a sum of $\KL$-divergences between
$\Bern(p_j^k)$ and $\Bern(r)$ over all candidate events in $\G^k$.

\paragraph{Total objective.}
\begin{equation}
\label{eq:total_loss}
\loss_{\text{total}} = \loss_{\text{cls}} + \beta\,\loss_{\text{MI}}.
\end{equation}

\subsection{Discrepancies in the Reference Implementation}

The official codebase\footnote{\url{https://github.com/sang-woo-seo/TGIB}}
implements the total loss as:
\begin{equation}
\label{eq:code_loss}
\loss_{\text{code}} = \BCE(p^+, 1) + \BCE(p^-, 0) + \loss_{\text{MI}},
\end{equation}
where $\BCE$ denotes binary cross-entropy, and $p^+$/$p^-$ are the
positive/negative link probabilities respectively. Comparing
\Cref{eq:total_loss} with~\Cref{eq:code_loss}, we identify the
following gaps:

\begin{enumerate}[leftmargin=*,itemsep=2pt]
  \item \textbf{Implicit $\beta = 1$:} the rate term enters the loss
        without any explicit weighting.
  \item \textbf{Non-standard $\KL$:} the implementation uses
        $\log(p/r + 10^{-6})$ instead of $\log p - \log r$.
  \item \textbf{Fixed temperature:} the Concrete sampler uses
        $\tau = 1$ throughout training with no annealing.
  \item \textbf{Hardcoded prior schedule:} $r$ anneals from $0.9$
        to $0.7$ with fixed hyperparameters.
  \item \textbf{No $\beta$-warmup:} $\beta$ is constant across all
        epochs.
\end{enumerate}

The remainder of this report addresses each of these in turn.

% ══════════════════════════════════════════════════════════════════════
\section{Improvement~1: Explicit $\beta$-Weighted Rate Term}
\label{sec:beta}
% ══════════════════════════════════════════════════════════════════════

\subsection{Motivation}

In the Information Bottleneck framework, $\beta$ governs the fundamental
tradeoff between \emph{compression} (small $\R^k$, i.e.\ sparse
explanations) and \emph{sufficiency} (high predictive accuracy). Setting
$\beta = 1$ fixes this tradeoff at a single point, forfeiting the
ability to navigate the rate--distortion curve.

For the \emph{explanation} task, one typically prefers \emph{sparser}
subgraphs than those optimised purely for prediction accuracy. An
explicit $\beta > 1$ increases the penalty for including
non-informative candidate events, yielding more selective explanations.

\subsection{Mathematical Formulation}

\begin{definition}[Explicit $\beta$-weighted objective]
\label{def:beta_objective}
For a hyperparameter $\beta > 0$, the improved TGIB objective is:
\begin{equation}
\label{eq:beta_objective}
\loss_{\text{total}}^{(\beta)} = \loss_{\text{cls}} + \beta \cdot \loss_{\text{MI}}.
\end{equation}
\end{definition}

When $\beta = 1$, this reduces identically to the current formulation.
The recommended operating region is $\beta \in [0.5, 5.0]$:
\begin{itemize}[leftmargin=*,itemsep=1pt]
  \item $\beta < 1$: relaxes compression; useful for maximising
        link-prediction accuracy at the cost of explanation density.
  \item $\beta > 1$: enforces compression; useful for generating
        sparse, interpretable explanations.
\end{itemize}

\subsection{Implementation Sketch}

\begin{algorithm}[H]
\caption{Forward pass with explicit $\beta$}
\label{alg:beta}
\begin{algorithmic}[1]
\REQUIRE Target mask scores $p^+, p^-$; prior $r$; weight $\beta > 0$
\STATE Compute $\loss_{\text{MI}}$ via \Cref{eq:mi_stable}
\STATE $\loss_{\text{MI}} \leftarrow \beta \cdot \loss_{\text{MI}}$
\RETURN $\loss_{\text{cls}} + \loss_{\text{MI}}$
\end{algorithmic}
\end{algorithm}

% ══════════════════════════════════════════════════════════════════════
\section{Improvement~2: Numerically Stable KL Divergence}
\label{sec:kl_stability}
% ══════════════════════════════════════════════════════════════════════

\subsection{The Current Implementation}

The reference code computes the per-event rate term as:
\begin{equation}
\label{eq:kl_code}
\ell_j^{\text{code}}
  = p_j \,\log\!\left(\frac{p_j}{r} + \epsilon\right)
  + (1 - p_j)\,\log\!\left(\frac{1 - p_j}{1 - r} + \epsilon\right),
  \qquad \epsilon = 10^{-6}.
\end{equation}

\subsection{Analysis of the Epsilon Placement}

This formulation has two problematic regimes:

\paragraph{Case 1: $p_j \to 0$.}
The first term becomes $0 \cdot \log(10^{-6}) = 0$, which is
correct. However, the gradient with respect to $p_j$ receives a
contribution of
$\log(\epsilon) - 1 \approx -15.8$, which is disproportionately large
relative to the true gradient $\log(p_j/r)$. This introduces spurious
gradient magnitude early in training when many mask probabilities are
near zero.

\paragraph{Case 2: $r \to 1$.}
The second term's ratio $(1 - p_j)/(1 - r)$ diverges, and the
epsilon-inside-log form clips this to $\log(1/\epsilon) \approx 13.8$,
artificially bounding the KL value and distorting the gradient.

\subsection{The Stable Formulation}

The $\KL$-divergence between two Bernoulli distributions admits a clean
closed form that can be computed in log-space:

\begin{proposition}[Stable $\KL(\Bern(p) \| \Bern(r))$]
\label{prop:kl_stable}
For $p \in [0,1]$ and $r \in (0,1)$, the $\KL$-divergence is:
\begin{equation}
\label{eq:kl_stable}
\KL\!\big(\Bern(p) \,\|\, \Bern(r)\big)
  = p\,\big(\log(p + \delta) - \log(r + \delta)\big)
  + (1 - p)\,\big(\log(1 - p + \delta) - \log(1 - r + \delta)\big),
\end{equation}
where $\delta > 0$ is a machine-precision constant (e.g.\
$\delta = 10^{-8}$).
\end{proposition}

\begin{proof}
Expanding the definition:
\begin{align}
\KL(\Bern(p) \| \Bern(r))
  &= p\,\log\frac{p}{r} + (1 - p)\,\log\frac{1 - p}{1 - r}
    \notag\\
  &= p\,(\log p - \log r) + (1 - p)\,(\log(1 - p) - \log(1 - r)).
    \label{eq:kl_expanded}
\end{align}
Adding $\delta$ inside each $\log$ argument avoids $\log(0)$ while
preserving the identity for $p, r \in (0, 1)$, since
$\log(x + \delta) = \log x + \log(1 + \delta/x)$ and
$\delta/x \ll 1$ for the typical operating range. The gradient of
\Cref{eq:kl_stable} with respect to~$p$ is well-conditioned:
\begin{equation}
\frac{\partial\, \KL}{\partial\, p}
  = \log(p + \delta) - \log(r + \delta)
    - \log(1 - p + \delta) + \log(1 - r + \delta),
\end{equation}
which has bounded magnitude for all $p \in [0, 1]$ and $r \in (0, 1)$.
\end{proof}

\subsection{Comparison with the Code Form}

The difference between \Cref{eq:kl_code} and \Cref{eq:kl_stable} is
the substitution:
\begin{equation}
\log\!\left(\frac{a}{b} + \epsilon\right)
  \;\longrightarrow\;
\log(a + \delta) - \log(b + \delta).
\end{equation}
These agree when $\epsilon \ll a/b$ but diverge when $a \ll b \cdot
\epsilon$. The log-space form is the standard pattern used in
numerical libraries (e.g.\
\texttt{torch.nn.functional.kl\_div}).

\subsection{Full Rate Loss with Stable KL}

Summing over all candidate events in $\G^k$:
\begin{equation}
\label{eq:mi_stable}
\loss_{\text{MI}}^{(\text{stable})}
  = \frac{1}{|\G^k|}\sum_{e_j \in \G^k}
    \KL\!\big(\Bern(p_j^k) \,\|\, \Bern(r)\big),
\end{equation}
where each $\KL$ is computed via \Cref{eq:kl_stable}. The
$\frac{1}{|\G^k|}$ normalisation ensures scale invariance across
different neighbourhood sizes.

% ══════════════════════════════════════════════════════════════════════
\section{Improvement~3: Concrete Temperature Annealing}
\label{sec:temp_anneal}
% ══════════════════════════════════════════════════════════════════════

\subsection{Background on the Concrete Distribution}

The Gumbel--Softmax (Concrete) distribution~\cite{jang2016gumbel}
provides a differentiable relaxation of discrete Bernoulli sampling.
For a Bernoulli variable with log-odds $\log p_j^k$, the
reparameterised sample is:
\begin{equation}
\label{eq:concrete}
\alpha_j^k
  = \sigma\!\left(
    \frac{\log p_j^k - \log(1 - p_j^k) + g_1 - g_2}{\tau}
  \right),
  \qquad g_1, g_2 \sim \text{Gumbel}(0, 1),
\end{equation}
where $\tau > 0$ is the \emph{temperature} parameter. As
$\tau \to 0^+$, the Concrete distribution converges to the discrete
Bernoulli distribution; as $\tau \to \infty$, samples become uniformly
distributed.

\subsection{The Fixed-$\tau$ Limitation}

The current implementation fixes $\tau = 1$ throughout training. At
$\tau = 1$, the Concrete distribution has significant variance: the
sampled masks are smooth realisations in $(0, 1)$ rather than
near-binary $\{0, 1\}$. This has two consequences:

\begin{enumerate}[leftmargin=*,itemsep=2pt]
  \item \textbf{Explanation quality:} the importance scores $p_j^k$
        never become cleanly binary, so the extracted explanation
        subgraph is inherently ambiguous.
  \item \textbf{Rate--distortion gap:} the $\KL$-divergence in
        \Cref{eq:kl_stable} is computed on the \emph{sampled} mask
        values, which are smoothed by the high temperature. The true
        $\KL$ of the underlying Bernoulli is not what is optimised.
\end{enumerate}

\subsection{Temperature Annealing Schedule}

Following standard practice in Gumbel--Softmax
training~\cite{maddison2017concrete,jang2016gumbel}, we propose an
exponential decay schedule for the temperature:

\begin{definition}[Exponential temperature schedule]
\label{def:temp_schedule}
Given an initial temperature $\tau_0 > 0$, a minimum temperature
$\tau_{\min} > 0$, and a decay rate $\rho \in (0, 1)$, the temperature
at epoch~$e$ is:
\begin{equation}
\label{eq:temp_schedule}
\tau(e)
  = \max\!\big(\tau_{\min},\; \tau_0 \cdot \rho^{\,e-1}\big).
\end{equation}
\end{definition}

\begin{remark}[Recommended hyperparameters]
Typical values are $\tau_0 \in [1, 10]$, $\tau_{\min} = 0.1$, and
$\rho \in [0.999, 0.9999]$. At $\tau_{\min} = 0.1$, the Concrete
distribution is concentrated within $\pm 0.05$ of the underlying
Bernoulli parameter, yielding near-binary masks.
\end{remark}

The modified Concrete sampler becomes:
\begin{equation}
\label{eq:concrete_annealed}
\alpha_j^{k,(e)}
  = \sigma\!\left(
    \frac{\log p_j^k - \log(1 - p_j^k) + g_1 - g_2}{\tau(e)}
  \right).
\end{equation}

\subsection{Effect on Gradients}

The gradient of the Concrete log-density with respect to the log-odds
$\log p_j^k$ scales as $\mathcal{O}(1/\tau)$. A high initial
temperature therefore facilitates early exploration (large gradient
signals across all mask values), while a low final temperature focuses
learning on the most critical events (gradients concentrate near the
decision boundary $p_j^k = 0.5$).

\subsection{Optional: Hard-Concrete Extension}

For truly binary explanations, we recommend the
\emph{hard-Concrete} relaxation of Louizos et~al.~\cite{louizos2018deep}:
\begin{equation}
\label{eq:hard_concrete}
\tilde{\alpha}_j^k
  = \mathrm{clamp}\!\big(
    \alpha_j^k \cdot (b - a) + a,\; 0,\; 1
  \big),
  \qquad a = -0.1,\; b = 1.1,
\end{equation}
where $\alpha_j^k$ is the soft Concrete sample from
\Cref{eq:concrete_annealed} and $\mathrm{clamp}(\cdot, l, u)$ clips
to $[l, u]$. The stretch parameters $a = -0.1$ and $b = 1.1$ allow
the distribution to place nonzero mass at exactly $0$ and $1$, while
remaining continuous. The IB rate term $\loss_{\text{MI}}$ still
operates on the underlying probabilities $p_j^k$, not the
hard-concretised~$\tilde{\alpha}_j^k$.

% ══════════════════════════════════════════════════════════════════════
\section{Improvement~4: Configurable Prior Schedule}
\label{sec:prior_schedule}
% ══════════════════════════════════════════════════════════════════════

\subsection{Role of the Prior $r$}

The prior $r$ in \Cref{eq:rate} determines the ``target'' mask
probability: the $\KL$-divergence penalises deviations of $p_j^k$
from~$r$. A high~$r$ encourages the mask to be \emph{on} (i.e.\ to
include candidate events); a low~$r$ encourages the mask to be
\emph{off} (i.e.\ to exclude them).

In the reference code, $r$ is annealed via a step function:
\begin{equation}
\label{eq:r_step}
r(e) = \max\!\big(
  r_{\min},\;
  r_0 - \lfloor e / \Delta \rfloor \cdot \delta_r
\big),
\end{equation}
with $r_0 = 0.9$, $\Delta = 10$, $\delta_r = 0.1$, and $r_{\min} = 0.7$.
This means the final prior is $r = 0.7$, i.e.\ the model is pushed
toward including $\sim$70\% of candidate events in the explanation
graph.

\subsection{The Sparsity--Accuracy Tradeoff of $r$}

Since the expected mask value converges to~$r$ under the $\KL$
penalty, the fraction of retained candidate events is approximately~$r$.
This gives direct control over explanation sparsity:

\begin{table}[ht]
\centering
\caption{Effect of final prior $r_{\min}$ on explanation sparsity.}
\label{tab:prior_sparsity}
\begin{tabular}{@{}ccc@{}}
\toprule
$r_{\min}$ & Approx.\ retention rate & Use case \\
\midrule
0.3 & 30\% & Sparse, focused explanations \\
0.5 & 50\% & Balanced \\
0.7 & 70\% & Current default (code) \\
0.5 (fixed) & 50\% & Paper-equivalent (no annealing) \\
\bottomrule
\end{tabular}
\end{table}

\subsection{Generalised Schedule}

We generalise the step-function schedule to accept arbitrary
parameters:

\begin{definition}[Generalised prior schedule]
\label{def:prior_schedule}
For an initial prior $r_0$, final prior $r_{\min}$, decay interval
$\Delta$, and per-interval decrement $\delta_r$:
\begin{equation}
\label{eq:prior_general}
r(e) = \max\!\big(
  r_{\min},\;
  r_0 - \lfloor e / \Delta \rfloor \cdot \delta_r
\big).
\end{equation}
Setting $r_0 = r_{\min}$ and $\delta_r = 0$ disables annealing
(constant prior). Setting $\Delta = 1$ makes the decay linear rather
than stepwise.
\end{definition}

\begin{remark}[Paper fidelity]
The published paper specifies $r$ as ``a predefined
hyperparameter''~\cite{seo2024tgib}. To reproduce the paper's
conditions exactly, set $r_0 = r_{\min}$ and $\delta_r = 0$, yielding
a constant prior for all epochs. The step-function annealing in the
code is a training heuristic not present in the paper.
\end{remark}

% ══════════════════════════════════════════════════════════════════════
\section{Improvement~5: $\beta$-Warmup Strategy}
\label{sec:beta_warmup}
% ══════════════════════════════════════════════════════════════════════

\subsection{Motivation}

Early in training, the encoder parameters are randomly initialised, and
the affinity scores $p_j^k$ carry little signal. Applying a large
$\beta$ from the first epoch forces the rate term
$\beta \cdot \loss_{\text{MI}}$ to dominate the gradient, potentially
trapping the model in a degenerate state where all masks are pushed to
$r$ before the model has learned to distinguish informative from
non-informative events.

A \emph{warmup} phase allows the model to first learn meaningful
$p_j^k$ values (driven by $\loss_{\text{cls}}$ alone), before the
compression pressure is gradually introduced.

\subsection{Linear Warmup Schedule}

\begin{definition}[$\beta$-warmup]
\label{def:beta_warmup}
Given a target $\beta^* > 0$ and warmup duration $E_w \geq 0$, the
effective weight at epoch~$e$ is:
\begin{equation}
\label{eq:beta_warmup}
\beta(e) = \beta^* \cdot \min\!\left(
  1,\; \frac{e}{E_w}
\right).
\end{equation}
When $E_w = 0$, this reduces to $\beta(e) = \beta^*$ (no warmup).
\end{definition}

The resulting epoch-dependent total objective is:
\begin{equation}
\label{eq:total_warmup}
\loss_{\text{total}}(e) = \loss_{\text{cls}} + \beta(e) \cdot \loss_{\text{MI}}.
\end{equation}

\subsection{Practical Recommendations}

\begin{itemize}[leftmargin=*,itemsep=2pt]
  \item For $\beta^* \leq 2$: warmup is generally unnecessary; the
        rate term does not dominate.
  \item For $\beta^* > 2$: use $E_w \in [3, 5]$ epochs to allow the
        encoder to stabilise before compression.
  \item For very large $\beta^* > 5$: use $E_w \in [5, 10]$ and
        consider also annealing $\tau$ more slowly to maintain gradient
        flow.
\end{itemize}

% ══════════════════════════════════════════════════════════════════════
\section{Combined Objective: Summary}
\label{sec:combined}
% ══════════════════════════════════════════════════════════════════════

Collecting all improvements, the fully improved TGIB objective at
epoch~$e$ is:

\begin{equation}
\label{eq:final_objective}
\boxed{
\loss_{\text{total}}(e)
  = \underbrace{-\log\sigma\!\big(q_\theta(X_{e_k}, H_+^k)\big)
    - N \cdot \E_{\text{neg}}\!\big[
      \log\sigma\!\big(q_\theta(X_{e_k}^{\text{neg}}, H_-^k)\big)
    \big]}_{\displaystyle\loss_{\text{cls}}}
  + \;\beta(e)
    \cdot \underbrace{\frac{1}{|\G^k|} \sum_{e_j \in \G^k}
      \KL\!\big(\Bern(p_j^k) \,\|\, \Bern(r(e))\big)}_{\displaystyle\loss_{\text{MI}}^{(\text{stable})}},
}
\end{equation}
where:
\begin{align}
\label{eq:components}
\beta(e) &= \beta^* \cdot \min\!\left(1,\; \frac{e}{E_w}\right),
  \notag\\[4pt]
r(e) &= \max\!\big(
  r_{\min},\;
  r_0 - \lfloor e / \Delta \rfloor \cdot \delta_r
\big),
  \notag\\[4pt]
\tau(e) &= \max\!\big(
  \tau_{\min},\;
  \tau_0 \cdot \rho^{\,e-1}
\big),
  \notag\\[4pt]
\KL\!\big(\Bern(p) \| \Bern(r)\big)
  &= p\,\log\!\frac{p + \delta}{r + \delta}
  + (1 - p)\,\log\!\frac{1 - p + \delta}{1 - r + \delta},
  \notag\\[4pt]
\alpha_j^{k,(e)} &= \sigma\!\left(
  \frac{\log p_j^k - \log(1 - p_j^k) + g_1 - g_2}{\tau(e)}
\right).
\end{align}

\subsection{Hyperparameter Reference Table}

\begin{table}[ht]
\centering
\caption{Complete hyperparameter reference for the improved TGIB.}
\label{tab:hyperparams}
\begin{tabular}{@{}llll@{}}
\toprule
Symbol & Default & Range & Description \\
\midrule
$\beta^*$    & 1.0   & $[0.1, 10]$  & Target IB rate weight \\
$E_w$        & 0     & $[0, 10]$    & $\beta$-warmup epochs \\
$\tau_0$     & 5.0   & $[1, 10]$    & Initial Concrete temperature \\
$\tau_{\min}$& 0.1   & $[0.01, 0.5]$& Minimum temperature \\
$\rho$       & 0.9999& $(0, 1)$      & Temperature decay rate \\
$r_0$        & 0.9   & $(0, 1)$      & Initial prior probability \\
$r_{\min}$   & 0.7   & $(0, 1)$      & Final prior probability \\
$\Delta$     & 10    & $[1, 100]$   & Prior decay interval (epochs) \\
$\delta_r$   & 0.1   & $(0, 0.5)$   & Prior decay per interval \\
$\delta$     & $10^{-8}$ & ---       & KL numerical stability \\
\bottomrule
\end{tabular}
\end{table}

% ══════════════════════════════════════════════════════════════════════
\section{Implementation Pseudocode}
\label{sec:pseudocode}
% ══════════════════════════════════════════════════════════════════════

\begin{algorithm}[H]
\caption{Improved TGIB Forward Pass}
\label{alg:forward}
\begin{algorithmic}[1]
\REQUIRE Target event $e_k$, neighbourhood $\G^k$, epoch $e$
\REQUIRE Hyperparameters: $\beta^*$, $E_w$, $\tau_0$, $\tau_{\min}$,
         $\rho$, $r_0$, $r_{\min}$, $\Delta$, $\delta_r$, $\delta$
\ENSURE  Prediction $\hat{Y}_k$, loss $\loss_{\text{total}}$

\STATE \textbf{// Compute epoch-dependent scalars}
\STATE $\beta \leftarrow \beta^* \cdot \min(1,\; e / E_w)$ \hfill\COMMENT{C5: warmup}
\STATE $\tau \leftarrow \max(\tau_{\min},\; \tau_0 \cdot \rho^{e-1})$
         \hfill\COMMENT{C3: temp annealing}
\STATE $r \leftarrow \max(r_{\min},\; r_0 - \lfloor e / \Delta \rfloor
         \cdot \delta_r)$
         \hfill\COMMENT{C4: prior schedule}

\STATE \textbf{// Compute mask probabilities}
\FOR{each candidate $e_j \in \G^k$}
  \STATE $p_j^k \leftarrow \sigma\!\big(g(X_{e_k}, Z_{e_j})\big)$
\ENDFOR

\STATE \textbf{// Concrete sampling with annealed temperature}
\FOR{each $e_j \in \G^k$}
  \STATE Sample $g_1, g_2 \sim \text{Gumbel}(0,1)$
  \STATE $\alpha_j^k \leftarrow \sigma\!\big(
    (\log p_j^k - \log(1 - p_j^k) + g_1 - g_2) / \tau
  \big)$
  \STATE $\tilde{Z}_{e_j} \leftarrow \alpha_j^k \cdot Z_{e_j}$
         \hfill\COMMENT{Masked candidate embedding}
\ENDFOR

\STATE \textbf{// Readout and prediction}
\STATE $H_+^k \leftarrow \mathrm{Readout}\!\big(\{\tilde{Z}_{e_j}\}_{e_j \in \G^k}\big)$
\STATE $\hat{Y}_k \leftarrow \sigma\!\big(q_\theta(X_{e_k}, H_+^k)\big)$

\STATE \textbf{// Stable KL computation (C2)}
\STATE $\loss_{\text{MI}} \leftarrow 0$
\FOR{each $e_j \in \G^k$}
  \STATE $\loss_{\text{MI}} \mathrel{+}=
    p_j^k\,\big(\log(p_j^k + \delta) - \log(r + \delta)\big)$
    $\;+\; (1 - p_j^k)\,\big(\log(1 - p_j^k + \delta) - \log(1 - r + \delta)\big)$
\ENDFOR
\STATE $\loss_{\text{MI}} \leftarrow \loss_{\text{MI}} / |\G^k|$

\STATE \textbf{// Total loss with explicit beta (C1)}
\STATE $\loss_{\text{total}} \leftarrow \loss_{\text{cls}} + \beta \cdot \loss_{\text{MI}}$
         \hfill\COMMENT{C1: beta weight}

\RETURN $\hat{Y}_k$, $\loss_{\text{total}}$
\end{algorithmic}
\end{algorithm}

% ══════════════════════════════════════════════════════════════════════
\section{Expected Impact}
\label{sec:impact}
% ══════════════════════════════════════════════════════════════════════

We summarise the expected impact of each improvement across the two
evaluation axes of TGIB:

\begin{table}[ht]
\centering
\caption{Summary of expected improvements.}
\label{tab:impact}
\begin{tabular}{@{}lccc@{}}
\toprule
Improvement & Link pred.\ (AP) & Explanation sparsity & Training stability \\
\midrule
C1: Explicit $\beta$ & $\uparrow$ (tunable) & $\uparrow$ (tunable) & --- \\
C2: Stable KL & --- & --- & $\uparrow\uparrow$ \\
C3: Temp.\ annealing & $\uparrow$ & $\uparrow\uparrow$ & $\uparrow$ \\
C4: Config.\ prior & $\leftrightarrow$/$\uparrow$ & $\uparrow\uparrow$ & --- \\
C5: $\beta$-warmup & $\uparrow$ & $\uparrow$ & $\uparrow\uparrow$ \\
\bottomrule
\end{tabular}
\end{table}

\begin{itemize}[leftmargin=*,itemsep=2pt]
  \item \textbf{Link prediction (AP):} Improvements~C1, C3, and C5
        are expected to yield modest gains in AP by enabling a more
        favourable exploration of the rate--distortion tradeoff. The
        KL fix~(C2) primarily improves training stability, which
        indirectly benefits AP through better convergence.

  \item \textbf{Explanation sparsity:} Improvements~C3 (temperature
        annealing) and~C4 (configurable prior) have the largest impact.
        Temperature annealing produces near-binary masks that cleanly
        separate ``important'' from ``unimportant'' events. The prior
        $r_{\min}$ directly controls the fraction of retained events.

  \item \textbf{Training stability:} The KL fix~(C2) and
        $\beta$-warmup~(C5) jointly address the two main sources of
        training instability: spurious gradients at mask boundaries,
        and premature dominance of the rate term.
\end{itemize}

% ══════════════════════════════════════════════════════════════════════
\section{Conclusion}
\label{sec:conclusion}
% ══════════════════════════════════════════════════════════════════════

This report has presented five mathematically grounded improvements to
the TGIB framework~\cite{seo2024tgib}, each addressing a specific
limitation in the reference implementation. The improvements are:

\begin{enumerate}[leftmargin=*,itemsep=2pt]
  \item \textbf{Explicit $\beta$-weighting} of the rate term, enabling
        principled control of the accuracy--explainability tradeoff.
  \item \textbf{Numerically stable $\KL$-divergence} via log-space
        computation, eliminating gradient artefacts near $p_j^k = 0$.
  \item \textbf{Temperature annealing} for the Concrete sampler,
        progressively discretising mask variables for crisp binary
        explanations.
  \item \textbf{Configurable prior schedule}, exposing the full
        sparsity--accuracy tradeoff controlled by $r_{\min}$.
  \item \textbf{$\beta$-warmup}, allowing the encoder to learn
        meaningful representations before compression pressure is
        applied.
\end{enumerate}

All improvements are drop-in compatible with the existing codebase and
require no architectural modifications to the TGIB model. They are
derived directly from the published GIB objective
(Equation~4 of~\cite{seo2024tgib}) and established practices in the
deep learning literature (Concrete
distributions~\cite{maddison2017concrete}, hard-Concrete
relaxations~\cite{louizos2018deep}, and IB
warmup~\cite{alemi2016deep}).

% ══════════════════════════════════════════════════════════════════════
% References
% ══════════════════════════════════════════════════════════════════════
\bibliographystyle{plain}
\begin{thebibliography}{10}

\bibitem{alemi2016deep}
A.~A. Alemi, I.~Fischer, J.~V. Dillon, and K.~P. Murphy.
\newblock Deep variational information bottleneck.
\newblock \emph{arXiv preprint arXiv:1612.00410}, 2016.

\bibitem{jang2016gumbel}
E.~Jang, S.~Gu, and B.~Poole.
\newblock Categorical reparameterization with {Gumbel}-Softmax.
\newblock In \emph{Proc. ICLR}, 2017.

\bibitem{louizos2018deep}
C.~Louizos, M.~Welling, and D.~P. Kingma.
\newblock Learning sparse neural networks through $L_0$ regularization.
\newblock In \emph{Proc. ICLR}, 2018.

\bibitem{maddison2017concrete}
C.~J. Maddison, A.~Mnih, and Y.~W. Teh.
\newblock The {Concrete} distribution: A continuous relaxation of discrete
  random variables.
\newblock In \emph{Proc. ICML}, 2017.

\bibitem{seo2024tgib}
S.~Seo, S.~Kim, J.~Jung, Y.~Lee, and C.~Park.
\newblock Self-explainable temporal graph networks based on graph
  information bottleneck.
\newblock In \emph{Proc. KDD}, pages 2572--2583, 2024.

\bibitem{wu2020gib}
T.~Wu, H.~Ren, P.~Li, and J.~Leskovec.
\newblock Graph information bottleneck.
\newblock In \emph{Proc. NeurIPS}, 2020.

\end{thebibliography}

\end{document}
